\title{Compilation et Interprétation} % Matière

\newcommand{\name}{Chaya Aouate} % Prénom Nom
\newcommand{\studentid}{15EK774K} % n° étudiant
\newcommand{\subtitle}{Chapitre 2 - 5 Le langage Assembleur \\
Chapitre 3 : Comprendre un programme compilé en C \\
Chapitre 4 : L'analyse Lexicale et Syntaxique \\
Chapitre 5 : Utilisation du compilateur Yacc} % Titre du devoir

% la commande section donne maintenant: 
% \section{À rendre} => Exercice 1.2 À rendre
% à commenter pour remettre les titres de sections habituels
\titleformat{\section}{\bfseries \large}{Exercice \thesection}{1em}{}

% Définition du style principal des codes sources
% styles définis dans src/codeblocks.tex
\lstdefinestyle{CodeBox}{
    style=tab4, % la largeur de tabulation vaut 4 espaces
    style=linesnb5, % numéro de lignes, toutes les 5 lignes
    style=box, % enferme les codes sources dans des boites colorées
    style=ttstyle, % police type typewriter
    style=syntaxcol, % coloration syntaxique
}